%========================================
% MatinBook — Stage 01: Basic Infrastructure Test
% Version: v1.1
% Purpose:
%   Validate the foundational systems of MatinBook v1.1:
%     1. Class loading and metadata
%     2. Frontmatter / Mainmatter geometry split
%     3. Color system (CMYK palette, 5 families × 3 shades)
%     4. Font system (Persian, Latin, Math, Monospace, Code)
%     5. Mathematical operators and delimiters
%     6. Basic box families (matinbox + matin-outline)
%     7. Typography (microtype, line spacing, widow/orphan)
%     8. Cross-referencing (cleveref with Persian names)
%     9. Margin notes (RTL)
%    10. Full document structure (cover → back cover)
%
% Compilation:
%   xelatex -shell-escape stage01-basic.tex
%   xelatex -shell-escape stage01-basic.tex
%   (2 passes required for cover + cross-references)
%
% Expected outcome:
%   A professionally typeset PDF demonstrating that all core
%   systems are functional and visually consistent.
%========================================

%------------------------------------------------------------------------
% Search Path (for tests directory)
%------------------------------------------------------------------------
\makeatletter
\def\input@path{{../}}
\makeatother

%------------------------------------------------------------------------
% Document Class
%------------------------------------------------------------------------
\documentclass{matinbook}

%------------------------------------------------------------------------
% Book Metadata
%------------------------------------------------------------------------
\title{آزمون زیرساخت پایه}
\author{تیم توسعه MatinBook}
\date{مهر ۱۴۰۵}

\booktitle{آزمون زیرساخت پایه}

\renewcommand{\repository}{github.com/matinmahmoudi-cs/matinbook}

%========================================================================
% DOCUMENT
%========================================================================
\begin{document}

%------------------------------------------------------------------------
% FRONT COVER
%------------------------------------------------------------------------
\makecover
    {آزمون زیرساخت پایه}
    {اعتبارسنجی سیستم‌های بنیادین MatinBook}
    {تیم توسعه MatinBook}
    {مهر ۱۴۰۵}

%------------------------------------------------------------------------
% FRONT MATTER (symmetric margins)
%------------------------------------------------------------------------
\frontmatter
\frontmattergeometry

% Title page
\maketitle

% Copyright / Colophon
\thispagestyle{empty}
\vfill
\begin{center}
    {\large\textbf{شناسنامه آزمون}}

    \vspace{0.8cm}

    \begin{tabular}{rl}
        عنوان: & آزمون زیرساخت پایه \\
        نسخه: & \matinbookversion \\
        تاریخ: & \matinbookdate \\
        موتور: & \lr{XeLaTeX} \\
        مجوز: & \lr{MIT} \\
    \end{tabular}

    \vspace{0.8cm}

    {\small این سند بخشی از مجموعه آزمون‌های یکپارچگی \lr{MatinBook} است.}
\end{center}
\clearpage

% Preface
\chapter*{پیشگفتار}
\addcontentsline{toc}{chapter}{پیشگفتار}

این سند، نخستین آزمون از مجموعه‌ی آزمون‌های یکپارچگی \lr{MatinBook} نسخه‌ی
\lr{\matinbookversion} است. هدف آن، اعتبارسنجی سیستم‌های بنیادین کلاس
پیش از آزمودن ماژول‌های تخصصی است.

در این آزمون، ده زیرسیستم کلیدی بررسی می‌شوند: بارگذاری کلاس و فراداده،
تفکیک هندسه‌ی پیش‌متن و متن اصلی، سیستم رنگ \lr{CMYK}، سیستم فونت،
عملگرهای ریاضی، خانواده‌های جعبه، تایپوگرافی، ارجاع هوشمند، یادداشت‌های
حاشیه‌ای و ساختار کامل سند از جلد تا پشت‌جلد.

هر بخش با مثال‌های عملی و ارجاعات متقابل همراه است تا خواننده بتواند
درستی رفتار هر زیرسیستم را در بافت واقعی یک کتاب ارزیابی کند.

\clearpage

% Table of Contents
\tableofcontents
\clearpage

% List of Figures
\listoffigures
\clearpage

% List of Tables
\listoftables
\clearpage

%------------------------------------------------------------------------
% MAIN MATTER (wide margin for notes)
%------------------------------------------------------------------------
\mainmatter
\mainmattergeometry

%========================================================================
% CHAPTER 1 — COLOR SYSTEM
%========================================================================
\chapter{سیستم رنگ‌ها}
\label{chap:colors}

\section{فلسفه‌ی رنگ در MatinBook}
\label{sec:color-philosophy}

سیستم رنگ \lr{MatinBook} بر پایه‌ی سه اصل طراحی شده است: \textbf{وضوح}
بر تزئین، \textbf{معناداری} بر تنوع، و \textbf{سازگاری} بر نوآوری.
هر رنگ در این سیستم، نقش معنایی مشخصی دارد و به‌جای آنکه صرفاً
جنبه‌ی زیبایی‌شناختی داشته باشد، به خواننده در تشخیص سریع نوع محتوا
کمک می‌کند.

مدل رنگ به‌کاررفته در \lr{MatinBook}، مدل \lr{CMYK} است که مطابق
استاندارد نشر آموزشی ایران (ز/۱-۴) انتخاب شده است. تمام درصدها
مضرب ۵ یا ۱۰ هستند تا در چاپ حرفه‌ای، رنگ‌های یکنواخت و قابل
پیش‌بینی تولید شوند. رنگ غالب سیستم، \texttt{matinblue} است و
سایر رنگ‌ها یا سایه‌های آبی هستند یا رنگ‌های عملکردی.

\mnote{مدل \lr{CMYK} برای چاپ، و مدل \lr{RGB} برای نمایش دیجیتال مناسب است.}

\section{خانواده‌ی آبی — قضایا و لم‌ها}
\label{sec:color-blue}

خانواده‌ی آبی، رنگ غالب سیستم است و برای محیط‌های نظری مانند
قضیه، لم، نتیجه و گزاره به کار می‌رود. این انتخاب بر اساس
سنت نشر دانشگاهی است که رنگ آبی را به مفاهیم نظری و دقیق
اختصاص می‌دهد.

\begin{itemize}
    \item \textcolor{matinblue}{\textbf{آبی اصلی}} — \texttt{matinblue}،
          \lr{CMYK(80, 30, 0, 30)}
    \item \textcolor{matindarkblue}{\textbf{آبی تیره}} — \texttt{matindarkblue}،
          \lr{CMYK(80, 30, 0, 60)}
    \item \colorbox{matinlightblue}{\textbf{آبی روشن}} — \texttt{matinlightblue}،
          \lr{CMYK(15, 5, 0, 5)}
\end{itemize}

\section{خانواده‌ی سبز — تعاریف و راه‌حل‌ها}
\label{sec:color-green}

خانواده‌ی سبز به مفاهیم تعریفی و راه‌حل‌ها اختصاص دارد. رنگ سبز
در روان‌شناسی رنگ‌ها، حس آرامش و رشد را منتقل می‌کند که با
ماهیت تعریف‌ها — یعنی ساختن پایه‌های دانش — هم‌خوانی دارد.

\begin{itemize}
    \item \textcolor{matingreen}{\textbf{سبز اصلی}} — \texttt{matingreen}،
          \lr{CMYK(80, 0, 45, 30)}
    \item \textcolor{matindarkgreen}{\textbf{سبز تیره}} — \texttt{matindarkgreen}،
          \lr{CMYK(80, 0, 45, 65)}
    \item \colorbox{matinlightgreen}{\textbf{سبز روشن}} — \texttt{matinlightgreen}،
          \lr{CMYK(10, 0, 5, 5)}
\end{itemize}

\section{خانواده‌ی نارنجی — مثال‌ها}
\label{sec:color-orange}

خانواده‌ی نارنجی برای مثال‌ها به کار می‌رود. نارنجی رنگی
گرم و جلب‌کننده‌ی توجه است و به‌خوبی می‌تواند خواننده را به
سمت کاربردهای عملی مفاهیم هدایت کند.

\begin{itemize}
    \item \textcolor{matinorange}{\textbf{نارنجی اصلی}} — \texttt{matinorange}،
          \lr{CMYK(0, 45, 85, 10)}
    \item \textcolor{matindarkorange}{\textbf{نارنجی تیره}} — \texttt{matindarkorange}،
          \lr{CMYK(0, 45, 85, 55)}
    \item \colorbox{matinlightorange}{\textbf{نارنجی روشن}} — \texttt{matinlightorange}،
          \lr{CMYK(0, 5, 15, 0)}
\end{itemize}

\section{خانواده‌ی قرمز — نکته‌ها و هشدارها}
\label{sec:color-red}

خانواده‌ی قرمز به نکته‌ها و هشدارها اختصاص دارد. استفاده‌ی
سنجیده از قرمز، به‌دلیل قدرت بصری بالای آن، ضروری است؛
زیرا استفاده‌ی بیش از حد، اثر هشداردهندگی آن را از بین می‌برد.

\begin{itemize}
    \item \textcolor{matinred}{\textbf{قرمز اصلی}} — \texttt{matinred}،
          \lr{CMYK(0, 65, 75, 10)}
    \item \textcolor{matindarkred}{\textbf{قرمز تیره}} — \texttt{matindarkred}،
          \lr{CMYK(0, 65, 75, 55)}
    \item \colorbox{matinlightred}{\textbf{قرمز روشن}} — \texttt{matinlightred}،
          \lr{CMYK(0, 15, 15, 0)}
\end{itemize}

\section{خانواده‌ی بنفش — تمرین‌ها}
\label{sec:color-purple}

خانواده‌ی بنفش به تمرین‌ها اختصاص دارد. بنفش، رنگی است که
مرز میان مفاهیم نظری و کاربردهای عملی را نشان می‌دهد و
به همین دلیل برای تمرین‌ها — که پلی میان نظریه و عمل هستند —
انتخاب شده است.

\begin{itemize}
    \item \textcolor{matinpurple}{\textbf{بنفش اصلی}} — \texttt{matinpurple}،
          \lr{CMYK(20, 60, 0, 30)}
    \item \textcolor{matindarkpurple}{\textbf{بنفش تیره}} — \texttt{matindarkpurple}،
          \lr{CMYK(20, 60, 0, 65)}
    \item \colorbox{matinlightpurple}{\textbf{بنفش روشن}} — \texttt{matinlightpurple}،
          \lr{CMYK(5, 10, 0, 5)}
\end{itemize}

\section{رنگ‌های خنثی}
\label{sec:color-neutral}

رنگ‌های خنثی، ستون فقرات تایپوگرافی \lr{MatinBook} هستند.
متن اصلی نه با سیاه مطلق، بلکه با خاکستری تیره‌ی نزدیک به سیاه
حروف‌چینی می‌شود تا کنتراست چشم‌نوازتری ایجاد شود و
خستگی چشم خواننده در مطالعه‌ی طولانی کاهش یابد.

\begin{itemize}
    \item \textcolor{matindarkgray}{\textbf{خاکستری تیره}} — متن اصلی
    \item \textcolor{matinmediumgray}{\textbf{خاکستری متوسط}} — متن فرعی
    \item \colorbox{matinlightgray}{\textbf{خاکستری روشن}} — پس‌زمینه
    \item \colorbox{matinbordergray}{\textbf{خاکستری مرزی}} — حاشیه‌ها
    \item \colorbox{matincream}{\textbf{کِرِم}} — پس‌زمینه‌ی گرم
\end{itemize}

\section{نمایش گرافیکی پالت}
\label{sec:color-swatches}

نمودار \cref{fig:color-swatches} پالت رنگی \lr{MatinBook} را به‌صورت
بصری نمایش می‌دهد. هر خانواده با سه سایه — اصلی، روشن و تیره —
نمایش داده شده است.

\begin{figure}[h]
\centering
\begin{latin}
\begin{tikzpicture}[scale=0.9]
    \foreach \color [count=\i] in {
        matinblue, matingreen, matinorange, matinred, matinpurple} {
        \fill[\color] (\i*2.4-1.5, 2.4) rectangle (\i*2.4-0.4, 3.2);
        \fill[\color!60] (\i*2.4-1.5, 1.4) rectangle (\i*2.4-0.4, 2.2);
        \fill[\color!25] (\i*2.4-1.5, 0.4) rectangle (\i*2.4-0.4, 1.2);
        \node[font=\tiny, anchor=north] at (\i*2.4-0.95, 0.2) {\color};
    }
    \draw[matindarkgray, thick] (-0.2, 0.2) -- (11.5, 0.2);
    \draw[matindarkgray, thick] (-0.2, 0.2) -- (-0.2, 3.4);
    \node[font=\footnotesize, anchor=west] at (0.1, 3.5) {سایه‌ها};
\end{tikzpicture}
\end{latin}
\caption{پالت رنگی \lr{MatinBook} — پنج خانواده، هر یک با سه سایه}
\label{fig:color-swatches}
\end{figure}

همان‌طور که در \cref{fig:color-swatches} پیداست، هر خانواده
دارای سه سایه است که امکان استفاده‌ی معنادار در سطوح مختلف
محتوا را فراهم می‌کند.

%========================================================================
% CHAPTER 2 — FONT SYSTEM
%========================================================================
\chapter{سیستم فونت‌ها}
\label{chap:fonts}

\section{فونت متن فارسی}
\label{sec:font-persian}

فونت پیش‌فرض متن فارسی در \lr{MatinBook}، \lr{XB Niloofar} است
که با مقیاس ۱.۲ بارگذاری می‌شود. این فونت به‌دلیل پوشش کامل
حروف فارسی و خوانایی بالا در متون فنی انتخاب شده است.

این یک پاراگراف طولانی فارسی است که به‌عنوان نمونه‌ی متن
برای ارزیابی فونت به کار می‌رود. هدف آن است که رفتار فونت
در مواجهه با حروف خاص فارسی مانند «گ»، «چ»، «پ» و «ژ» بررسی شود.
همچنین ترکیب این حروف با حروف دیگر در کلمات طولانی مانند
«برنامه‌نویسی»، «الگوریتم‌های» و «پیاده‌سازی» آزموده می‌شود.

اعداد فارسی: ۰ ۱ ۲ ۳ ۴ ۵ ۶ ۷ ۸ ۹

اعداد لاتین در متن فارسی: \lr{0 1 2 3 4 5 6 7 8 9}

نشانه‌های نگارشی فارسی: «نقل قول»، (پرانتز)، [کروشه]، \{آکولاد\}،
نقطه‌ویرگول؛ و ویرگول،

\section{فونت متن لاتین}
\label{sec:font-latin}

\begin{latin}
The default Latin font in MatinBook is Times New Roman, with
TeX Gyre Termes as a fallback. The font is chosen for its
professional appearance and excellent readability in technical
documents.

The quick brown fox jumps over the lazy dog.
Pack my box with five dozen liquor jugs.

Special ligatures: fi fl ff ffi ffl

Quotation marks: ``double'' and `single'

Mathematical symbols in text: \(\alpha, \beta, \gamma, \delta, \epsilon\)
\end{latin}

\section{ترکیب فارسی و لاتین}
\label{sec:font-mixing}

در متن‌های فنی فارسی، ترکیب متن فارسی و لاتین امری اجتناب‌ناپذیر
است. در \lr{MatinBook}، تمام واژه‌ها و عبارات لاتین باید در
دستور \lr{\textbackslash lr\{...\}} قرار گیرند تا جهت‌نمایی
متن به‌درستی حفظ شود.

این یک متن فارسی است که شامل واژه‌های لاتین مانند \lr{computer}،
\lr{software}، \lr{algorithm}، \lr{function} و \lr{variable} است.
همچنین اعداد لاتین مانند \lr{3.14159} و \lr{42} در متن فارسی
قرار گرفته‌اند.

نام‌های خاص خارجی: \lr{Donald Knuth}، \lr{Leslie Lamport}،
\lr{Alan Turing}، \lr{John von Neumann}.

\section{فونت ریاضی}
\label{sec:font-math}

فونت ریاضی پیش‌فرض \lr{MatinBook}، \lr{XITS Math} است که
پوشش گسترده‌ای از نمادهای ریاضی دارد. در ادامه، چند نمونه
از کاربرد این فونت نمایش داده می‌شود.

معادله‌ی درون‌خطی: $E = mc^2$؛ قضیه‌ی فیثاغورث: $a^2 + b^2 = c^2$.

معادله‌ی نمایشی:
\begin{equation}
    \int_{0}^{\infty} \frac{x^3}{e^x - 1} \,dx = \frac{\pi^4}{15}
    \label{eq:stefan-boltzmann}
\end{equation}

معادله‌ی \cref{eq:stefan-boltzmann} به‌عنوان یک نمونه‌ی کلاسیک
از انتگرال‌های مهم در فیزیک آماری شناخته می‌شود.

حروف یونانی: $\alpha, \beta, \gamma, \Gamma, \delta, \Delta,
\epsilon, \varepsilon, \zeta, \eta, \theta, \Theta, \lambda,
\Lambda, \mu, \pi, \Pi, \sigma, \Sigma, \tau, \phi, \Phi,
\omega, \Omega$.

عملگرهای بزرگ:
\[
    \sum_{n=1}^{\infty} \frac{1}{n^2} = \frac{\pi^2}{6}, \qquad
    \prod_{k=1}^{n} k = n!, \qquad
    \lim_{x \to \infty} \left(1 + \frac{1}{x}\right)^x = e
\]

\section{فونت مونواسپیس و کد}
\label{sec:font-mono}

برای کدهای درون‌خطی و بلوک‌های کد، \lr{MatinBook} از فونت‌های
\lr{DejaVu Sans Mono} و \lr{JetBrains Mono} استفاده می‌کند.
نمونه‌ی کد درون‌خطی: \inlcode{print("Hello, World!")}.

%========================================================================
% CHAPTER 3 — MATHEMATICAL SYSTEM
%========================================================================
\chapter{سیستم ریاضی}
\label{chap:math}

\section{عملگرهای ریاضی سفارشی}
\label{sec:math-operators}

\lr{MatinBook} مجموعه‌ای از عملگرهای ریاضی سفارشی را ارائه می‌دهد
که برای حوزه‌های مختلف ریاضیات و علوم کامپیوتر طراحی شده‌اند.

\subsection{حساب برداری}

عملگرهای گرادیان، دیورژانس و کرل:
\[
    \grad f = \nabla f, \qquad
    \diver \vec{F} = \nabla \cdot \vec{F}, \qquad
    \curl \vec{F} = \nabla \times \vec{F}
\]

\subsection{جبر خطی}

عملگرهای رتبه، رد و ماتریس قطری:
\[
    \rank(A) = 3, \qquad
    \trace(A) = \sum_{i=1}^{n} a_{ii}, \qquad
    \diag(a_1, a_2, \ldots, a_n)
\]

\subsection{بهینه‌سازی}

عملگرهای \lr{argmax} و \lr{argmin}:
\[
    \argmax_{x \in \R} f(x), \qquad \argmin_{x \in \R} g(x)
\]

\section{مجموعه‌های اعداد}
\label{sec:math-numbers}

مجموعه‌های اعداد استاندارد:
\[
    \N \subset \Z \subset \Q \subset \R \subset \C
\]

\section{جداکننده‌های هوشمند}
\label{sec:math-delimiters}

جداکننده‌های هوشمند که به‌طور خودکار با محتوای خود مقیاس می‌گیرند:

\begin{align}
    \abs{\frac{a}{b}} &= \frac{\abs{a}}{\abs{b}} \label{eq:abs} \\
    \norm{\sum_{i=1}^{n} x_i} &\leq \sum_{i=1}^{n} \norm{x_i} \label{eq:norm} \\
    \inner{u}{v} &= \sum_{i=1}^{n} u_i v_i \label{eq:inner}
\end{align}

معادلات \cref{eq:abs,eq:norm,eq:inner} نمونه‌هایی از
جداکننده‌های هوشمند را نشان می‌دهند.

\section{ماتریس‌ها}
\label{sec:math-matrices}

نمایش ماتریس‌ها با سه سبک مختلف:

\[
    A = \mat{1 & 2 & 3 \\ 4 & 5 & 6 \\ 7 & 8 & 9}, \qquad
    B = \pmat{1 & 0 \\ 0 & 1}, \qquad
    \det(C) = \vmat{1 & 2 \\ 3 & 4} = -2
\]

\section{مشتق و انتگرال}
\label{sec:math-derivatives}

نمایش مشتق و انتگرال:
\[
    \deriv{f}{x} = \lim_{h \to 0} \frac{f(x + h) - f(x)}{h}, \qquad
    \pderiv{f}{x} = \frac{\partial f}{\partial x}
\]

انتگرال معین:
\[
    \int_{a}^{b} f(x) \, \dd{x} = F(b) - F(a)
\]

%========================================================================
% CHAPTER 4 — BOX FAMILIES
%========================================================================
\chapter{خانواده‌های جعبه}
\label{chap:boxes}

\section{فلسفه‌ی دو خانواده}
\label{sec:box-philosophy}

\lr{MatinBook} دو خانواده‌ی جعبه ارائه می‌دهد که هر یک برای
نوع خاصی از محتوا طراحی شده‌اند. خانواده‌ی \lr{matinbox} با
نوار رنگی توپر و بدنه‌ی سفید، برای محتوای نظری و رسمی مناسب است.
خانواده‌ی \lr{matin-outline} با خط عمودی سمت راست و قواعد افقی،
برای محتوای توضیحی و کاربردی به کار می‌رود.

هر دو خانواده، قابلیت شکستن در چند صفحه را دارند؛ یعنی اگر
محتوای جعبه از یک صفحه فراتر رود، جعبه به‌طور خودکار در مرز
صفحه شکسته می‌شود و در صفحه‌ی بعد ادامه می‌یابد.

\section{خانواده‌ی matinbox (سبک Boyer)}
\label{sec:box-matinbox}

\subsection{قضیه}

\begin{theorem}[قضیه‌ی نمونه]
    این یک قضیه‌ی آزمایشی است که برای ارزیابی خانواده‌ی
    \lr{matinbox} نوشته شده است. عنوان جعبه با نوار رنگی
    توپر و متن سفید نمایش داده می‌شود.
    \label{thm:sample}
\end{theorem}

طبق \cref{thm:sample}، خانواده‌ی \lr{matinbox} به‌درستی کار می‌کند.

\subsection{لم، نتیجه و گزاره}

\begin{lemma}[لم نمونه]
    این یک لم آزمایشی است.
    \label{lem:sample}
\end{lemma}

\begin{corollary}[نتیجه‌ی نمونه]
    این یک نتیجه‌ی آزمایشی است.
    \label{cor:sample}
\end{corollary}

\begin{proposition}[گزاره‌ی نمونه]
    این یک گزاره‌ی آزمایشی است.
    \label{prop:sample}
\end{proposition}

قضیه‌ی \cref{thm:sample}، لم \cref{lem:sample}، نتیجه‌ی
\cref{cor:sample} و گزاره‌ی \cref{prop:sample} همگی از
یک شمارنده‌ی مشترک استفاده می‌کنند.

\subsection{تعریف، نکته و تمرین}

\begin{definition}[تعریف نمونه]
    این یک تعریف آزمایشی است.
    \label{def:sample}
\end{definition}

\begin{remark}
    این یک نکته‌ی آزمایشی است.
    \label{rem:sample}
\end{remark}

\begin{exercise}
    این یک تمرین آزمایشی است.
    \label{exc:sample}
\end{exercise}

\section{خانواده‌ی matin-outline (سبک RTL)}
\label{sec:box-outline}

\subsection{مثال}

\begin{example}[مثال نمونه]
    این یک مثال آزمایشی است که برای ارزیابی خانواده‌ی
    \lr{matin-outline} نوشته شده است. خط عمودی در سمت
    راست جعبه قرار دارد و عنوان با متن مشکی و بدون
    پس‌زمینه نمایش داده می‌شود.
    \label{ex:sample}
\end{example}

\subsection{اثبات}

\begin{proof}
    این یک اثبات آزمایشی است. مربع \lr{QED} در انتهای
    جعبه به‌طور خودکار اضافه می‌شود.
\end{proof}

\subsection{راه‌حل}

\begin{solution}
    این یک راه‌حل آزمایشی است.
\end{solution}

\section{شمارنده‌ی مشترک}
\label{sec:box-counter}

تمام محیط‌های شماره‌دار در \lr{MatinBook} از یک شمارنده‌ی
مشترک استفاده می‌کنند. این طراحی، بر اساس نظر جامعه‌ی
\lr{LaTeX} فارسی، به خواننده کمک می‌کند تا مفاهیم را
سریع‌تر پیدا کند. برای نمونه، اگر تعریفی با شماره‌ی ۱.۱
ثبت شود و سپس قضیه‌ای با شماره‌ی ۱.۲ بیاید، خواننده
می‌داند که این دو، در یک فصل و به‌ترتیب ظاهر شده‌اند.

%========================================================================
% CHAPTER 5 — TYPOGRAPHY
%========================================================================
\chapter{تایپوگرافی}
\label{chap:typography}

\section{فاصله‌ی خطوط}
\label{sec:typography-leading}

فاصله‌ی خطوط در \lr{MatinBook} روی ۱.۱۵ تنظیم شده است که
مطابق سبک کتاب‌های درسی دانشگاهی مانند \lr{Boyer} و
\lr{Stewart} است. این مقدار، تعادلی میان خوانایی و
فشردگی متن ایجاد می‌کند.

این یک پاراگراف طولانی فارسی است که برای ارزیابی
تایپوگرافی نوشته شده است. هدف آن است که ببینیم آیا
خطوط به‌خوبی تراز می‌شوند و فاصله‌ی میان کلمات مناسب
است یا خیر. تایپوگرافی خوب، خوانایی را افزایش می‌دهد
و خستگی چشم خواننده را کاهش می‌دهد.

\lr{MatinBook} از \lr{microtype} برای بهبود تایپوگرافی
استفاده می‌کند. این بسته، با اعمال \lr{protrusion} روی
نشانه‌های نگارشی، ترازبندی لبه‌ها را بهبود می‌بخشد.

\section{جلوگیری از خطوط آویزان و یتیم}
\label{sec:typography-widow}

در حروف‌چینی حرفه‌ای، باید از ایجاد خطوط آویزان (\lr{widow})
و یتیم (\lr{orphan}) جلوگیری شود. \lr{MatinBook} با تنظیم
جریمه‌های مناسب، این مشکلات را به حداقل می‌رساند.

\begin{definition}[خط آویزان]
    خط آویزان به آخرین خط یک پاراگراف گفته می‌شود که به
    تنهایی در بالای صفحه‌ی بعد قرار گیرد.
    \label{def:widow}
\end{definition}

\begin{definition}[خط یتیم]
    خط یتیم به اولین خط یک پاراگراف گفته می‌شود که
    در انتهای صفحه باقی بماند.
    \label{def:orphan}
\end{definition}

\lr{MatinBook} با تنظیم \lr{\textbackslash widowpenalty=10000}
و \lr{\textbackslash clubpenalty=10000} از این مشکلات
جلوگیری می‌کند.

\section{یادداشت‌های حاشیه‌ای}
\label{sec:typography-margin}

یکی از ویژگی‌های جدید \lr{MatinBook} نسخه‌ی \lr{\matinbookversion}،
پشتیبانی از یادداشت‌های حاشیه‌ای است. حاشیه‌ی بیرونی پهن
(۵.۹ سانتی‌متر) فضای کافی برای یادداشت‌های توضیحی فراهم می‌کند.

\mnote{یادداشت‌های حاشیه‌ای باید کوتاه و مکمل متن اصلی باشند.}

دستور \lr{\textbackslash mnote\{...\}} برای یادداشت‌های متنی
و دستور \lr{\textbackslash marginfig\{...\}\{...\}} برای
شکل‌های حاشیه‌ای به کار می‌رود.

%========================================================================
% CHAPTER 6 — CROSS-REFERENCES
%========================================================================
\chapter{ارجاع هوشمند}
\label{chap:cref}

\section{سیستم ارجاع}
\label{sec:cref-system}

\lr{MatinBook} از بسته‌ی \lr{cleveref} برای ارجاع هوشمند
استفاده می‌کند. با این سیستم، نیازی به نوشتن دستی نوع
مرجع (مانند «قضیه»، «شکل» یا «معادله») نیست؛ زیرا
\lr{cleveref} نوع مرجع را به‌طور خودکار تشخیص می‌دهد.

\section{نمونه‌های ارجاع}
\label{sec:cref-examples}

ارجاع به معادله: \cref{eq:stefan-boltzmann}

ارجاع به شکل: \cref{fig:color-swatches}

ارجاع به قضیه: \cref{thm:sample}

ارجاع به تعریف: \cref{def:sample}

ارجاع به فصل: \cref{chap:colors}

ارجاع به بخش: \cref{sec:color-blue}

ارجاع به چند مرجع: \cref{thm:sample,lem:sample,cor:sample}

ارجاع بازه‌ای: \crefrange{eq:abs}{eq:inner}

همان‌طور که در \cref{fig:color-swatches} دیده می‌شود،
سیستم ارجاع هوشمند \lr{MatinBook} به‌درستی کار می‌کند.

%========================================================================
% CHAPTER 7 — SUMMARY
%========================================================================
\chapter{خلاصه‌ی نتایج}
\label{chap:summary}

\section{وضعیت زیرسیستم‌ها}
\label{sec:summary-status}

\begin{itemize}
    \item \textcolor{matingreen}{\textbf{بارگذاری کلاس:}} موفق
    \item \textcolor{matingreen}{\textbf{فراداده:}} موفق
    \item \textcolor{matingreen}{\textbf{تفکیک هندسه:}} موفق
    \item \textcolor{matingreen}{\textbf{سیستم رنگ:}} ۲۰ رنگ \lr{CMYK} در ۵ خانواده
    \item \textcolor{matingreen}{\textbf{فونت فارسی:}} \lr{XB Niloofar} با مقیاس ۱.۲
    \item \textcolor{matingreen}{\textbf{فونت لاتین:}} \lr{Times New Roman}
    \item \textcolor{matingreen}{\textbf{فونت ریاضی:}} \lr{XITS Math}
    \item \textcolor{matingreen}{\textbf{فونت مونواسپیس:}} \lr{DejaVu Sans Mono}
    \item \textcolor{matingreen}{\textbf{عملگرهای ریاضی:}} ۱۴ عملگر سفارشی
    \item \textcolor{matingreen}{\textbf{جداکننده‌های هوشمند:}} ۶ جداکننده
    \item \textcolor{matingreen}{\textbf{خانواده‌های جعبه:}} \lr{matinbox} + \lr{matin-outline}
    \item \textcolor{matingreen}{\textbf{محیط‌های قضیه:}} ۱۱ محیط
    \item \textcolor{matingreen}{\textbf{تایپوگرافی:}} \lr{microtype} فعال
    \item \textcolor{matingreen}{\textbf{ارجاع هوشمند:}} \lr{cleveref} فارسی
    \item \textcolor{matingreen}{\textbf{یادداشت حاشیه:}} \lr{\textbackslash mnote}
\end{itemize}

\section{نتیجه‌گیری}
\label{sec:summary-conclusion}

\textbf{تمام زیرسیستم‌های پایه‌ی \lr{MatinBook} نسخه‌ی
\lr{\matinbookversion} با موفقیت بارگذاری و آزمایش شدند.
این آزمون، پایه‌ای برای آزمون‌های تخصصی‌تر فراهم می‌کند
و اطمینان می‌دهد که زیرساخت کلاس، آماده‌ی استفاده در
پروژه‌های واقعی است.}

%------------------------------------------------------------------------
% BACK MATTER (symmetric margins)
%------------------------------------------------------------------------
\backmatter
\frontmattergeometry

% Bibliography (empty placeholder)
\printbibliography[title={منابع و مراجع}]

% Index
\printindex

%------------------------------------------------------------------------
% BACK COVER
%------------------------------------------------------------------------
\makebackcover{%
    این سند، نخستین آزمون از مجموعه‌ی آزمون‌های یکپارچگی
    \lr{MatinBook} نسخه‌ی \lr{\matinbookversion} است.

    \vspace{0.5cm}

    \textbf{زیرسیستم‌های آزموده‌شده:}
    \begin{itemize}
        \item سیستم رنگ \lr{CMYK}
        \item فونت‌های فارسی، لاتین و ریاضی
        \item عملگرهای ریاضی سفارشی
        \item خانواده‌های جعبه
        \item تایپوگرافی و یادداشت حاشیه
        \item ارجاع هوشمند
    \end{itemize}
}

\end{document}